% D.60; full independent root proof audit passed.
\subsection{An infinitely ramified Alexander module with zero first cohomology}

\begin{statement}{D.60}{Proved}{A completed cusp tangent with a unique metabelian lift}
Let \(M=F^{(5)}/(F^{(5)})''\) be the geometrically pro-five
metabelian quotient of D.35, with its actual tangential arithmetic
action. There is an infinitely supported completed Kummer class
\(a\in H^1(\mathbb Q,\mathbb Z_5(1))\), with valuation one at two,
such that for the actual twisted Alexander module of D.35,
\[
          H^1_{\mathrm{cts}}(G_{\mathbb Q},A_{(a,0)})=0.
                                                        \tag{D.60.1}
\]
Consequently the abelian section class \((a,0)\) has exactly one
metabelian section class above it: the class of its completed cusp
tangent at zero. It is the image of an actual full cusp section.

The construction gives a genuine rank-zero example in D.47, whereas
D.44 gives rank \(\aleph_0\) for every abelian pair supported only
at five. It does not determine the rank for every arbitrarily
ramified pair, or the least derived length of a Fermat obstruction.
\end{statement}

\begin{proof}
Write \(\Lambda=\mathbb Z_5[[u,v]]\),
\(\mathfrak m=(5,u,v)\), and \(t_1=1+u\).
Choose once and for all compatible roots defining continuous
integral Kummer cocycles \(\kappa(q)\) for each rational prime
\(q\). We shall choose pairwise disjoint finite sets of primes
\(Q_r\), for \(r\ge1\), disjoint from \(\{2,5\}\), and put
\[
 a=\kappa(2)+\sum_{r\ge1}5^r\sum_{q\in Q_r}\kappa(q).
                                                        \tag{D.60.2}
\]
Each inner sum is finite. The outer series converges uniformly
as a series of \(\mathbb Z_5(1)\)-valued continuous cocycles:
the entire \(r\)-th summand takes values in
\(5^r\mathbb Z_5(1)\), irrespective of the number of terms.
Thus (D.60.2) is an actual continuous cocycle, not only a formal
family of local valuations. Distinct prime supports give its
valuation one at two and \(5^r\) at each prime of \(Q_r\).

\paragraph{Finite construction at stage \(r\).}
Suppose \(Q_1,\ldots,Q_{r-1}\) have been chosen. Set
\[
 a_{<r}=\kappa(2)+
       \sum_{j<r}5^j\sum_{q\in Q_j}\kappa(q),\qquad
 S_r=\{2,5,\infty\}\cup\bigcup_{j<r}Q_j .
                                                        \tag{D.60.3}
\]
Give the finite additive module
\(W_r=\Lambda C/\mathfrak m^r\Lambda C\)
the actual action \(\rho_r\) of the abelian pair
\((a_{<r},0)\). This action factors through
\(G_{\mathbb Q,S_r}\). Indeed the untwisted tripod action is
unramified away from five, and the finite Kummer sum is
unramified outside the displayed set.

Any later additions in (D.60.2) leave this action unchanged.
By (D.47.4), for every \(c\in\mathbb Z_5\),
\[
           t_1^{5^r c}-1\in\mathfrak m^{r+1}.
                                                        \tag{D.60.4}
\]
Formula (D.35.1) therefore makes \(\rho_r\) exactly the
reduction of the final action, for every possible continuation
of our construction.

The set
\[
             E_r=Z^1(G_{\mathbb Q,S_r},W_r)
                                                        \tag{D.60.5}
\]
of actual continuous cocycles is finite. This is the finite
coefficient, finite ramification finiteness used in D.47:
\(H^1(G_{\mathbb Q,S_r},W_r)\) and the group of coboundaries
are both finite. Equivalently the corresponding finite affine
representations are bounded-degree extensions unramified
outside the fixed finite set.

For each \(z\in E_r\) whose restriction to
\(\ker\rho_r\) is nonzero, choose
\(\delta\in\ker\rho_r\) with \(z(\delta)\ne0\).
The finite affine representation
\[
 g\longmapsto\bigl(z(g),\rho_r(g)\bigr)
       \in W_r\rtimes\operatorname{Aut}(W_r)
                                                        \tag{D.60.6}
\]
contains the element \((z(\delta),1)\).
Chebotarev \cite[Chapter V, Theorem 3.23]{PBC_MilneCFT}
supplies a prime \(q\notin S_r\) whose Frobenius
belongs to its conjugacy class. Choose these primes distinct
for the finitely many cocycles under consideration. Then
\[
                 \rho_r(\operatorname{Frob}_q)=1,
                 \qquad z(\operatorname{Frob}_q)\ne0.
                                                        \tag{D.60.7}
\]
The second property is independent of the Frobenius
representative: conjugating \((w,1)\) in the affine group
sends \(w\) to an invertible linear transform of \(w\).
Let \(Q_r\) contain all the chosen witness primes and one
additional unused prime different from two and five. The
additional prime ensures infinite support, including when
there are no cocycles to eliminate. This completes the
induction and defines (D.60.2).

\paragraph{Every unrestricted integral cocycle is captured.}
Let \(z:G_{\mathbb Q}\to A_{(a,0)}\) be any continuous
integral cocycle. Its reduction \(z_r\) belongs to the
particular finite set (D.60.5), even if \(z\) has unrestricted
ramification. Here is the full uniform support argument.
At a prime \(q\ne5\) with zero abelian inertia, D.47 proves
that the inertia value of \(z\) is zero: wild inertia and
the prime-to-five tame factors have no homomorphism to the
pro-five coefficient group, and the first nonzero augmentation
grade of a putative tame value would satisfy
\((q^{j+2}-q)m=0\) in a torsion-free lattice.
At a prime \(q\in Q_j\), D.43 gives
\(H^1(G_{\mathbb Q_q},A_{(a,0)})=0\), so the local cocycle
is an actual integral coboundary. Its tame inertia value is
\[
                        (t_1^{5^j}-1)h_q,
                      \qquad h_q\in A_{(a,0)}.
                                                        \tag{D.60.8}
\]
If \(j\ge r\), this lies in
\(\mathfrak m^{j+1}A_{(a,0)}\subset\mathfrak m^rA_{(a,0)}\).
The action has trivial inertia modulo \(\mathfrak m^r\) there
as well. These are all the primes outside \(S_r\).
Hence both the action and the cocycle graph are unramified
outside \(S_r\), which is precisely the assertion that
\(z_r\in E_r\). This argument does not presume that a
general finite coefficient cocycle lifts integrally, or
choose a ramification set depending on \(z\).

If \(z_r\) were nonzero on \(\ker\rho_r\), stage \(r\)
would have selected its witness prime \(q\in Q_r\).
The final inertia vector there is \((5^r,0)\ne(0,0)\),
so D.43 again expresses the restriction of \(z\) as an
integral coboundary, including on a Frobenius lift.
But (D.60.7) says that this Frobenius acts trivially on
\(W_r\). Reduction of that coboundary therefore gives
\(z_r(\operatorname{Frob}_q)=0\), contradicting
(D.60.7). We have proved
\[
                         z_r|_{\ker\rho_r}=0
                         \quad\hbox{for every }r.
                                                        \tag{D.60.9}
\]

\paragraph{Descent to the actual operator image.}
Let \(I=\rho_{(a,0)}(G_{\mathbb Q})\), with its compact
inverse-limit topology. An element of the kernel of this
full action belongs to the kernel of every \(\rho_r\).
By (D.60.9) its cocycle value is zero modulo every
\(\mathfrak m^r\), and hence is zero. The cocycle \(z\)
is consequently constant on fibres of
\(G_{\mathbb Q}\twoheadrightarrow I\) and descends to a
continuous cocycle on \(I\). Continuity follows from the
compact quotient topology; no passage from finite-level
cohomology through an unexamined derived limit occurs.

We show directly that \(H^1(I,A_{(a,0)})=0\).
The image of five-primary inertia at two is the subgroup
\[
 N=\{m\longmapsto t_1^c m:c\in\mathbb Z_5\}\subset I.
                                                        \tag{D.60.10}
\]
It is the full displayed subgroup because the valuation of
\(a\) at two is one. The original tripod action and
cyclotomic character are unramified there. Every operator
in (D.35.1) is a multiplier followed by the substitution
\(t_i\mapsto t_i^{\chi_5(g)}\). It conjugates multiplication
by \(t_1^c\) to multiplication by
\(t_1^{c\chi_5(g)}\); therefore \(N\) is normal in \(I\).
Moreover
\[
                         A_{(a,0)}^N=0,
                                                        \tag{D.60.11}
\]
since \(t_1-1=u\) is a nonzero divisor in \(\Lambda\).

Given a continuous cocycle \(h:I\to A_{(a,0)}\), pull it
back to \(G_{\mathbb Q_2}\). D.43, applied to the primitive
inertia vector \((1,0)\), makes this pullback a coboundary
of one element \(b\in A_{(a,0)}\). Subtract the same
coboundary on all of \(I\). The resulting cocycle \(h'\)
vanishes on \(N\). For \(n\in N\) and \(g\in I\), its
cocycle identity and normality of \(N\) give
\[
 n h'(g)=h'(ng)
        =h'\bigl(g(g^{-1}ng)\bigr)=h'(g).
\]
Thus \(h'(g)\in A_{(a,0)}^N=0\) for every \(g\).
This proves the desired vanishing on \(I\), and the descent
already proved gives (D.60.1).

Finally the cusp section \(g\mapsto x^{a(g)}s_0(g)\)
exists in the actual geometrically pro-five fundamental
group. It is also the pro-five image of a full cusp section:
in the product decomposition of \(\widehat{\mathbb Z}(1)\),
use \(a\) in the five-adic component and zero in the others.
Its abelian class is \((a,0)\). By D.35 the nonempty
metabelian fibre is a torsor under (D.60.1), and hence is
a singleton. In particular this conclusion concerns actual
section conjugacy classes, with the abelian conjugator
stabilizer already removed in D.35.
\end{proof}
